% !PRINTED 333
\weberpage{333}
%<<CONT
\[
\begin{aligned}
H' &= (-1)^{n-1}(n-1)\,N(\mathfrak{a})^{-n}(\lambda_1\lambda_2\ldots\lambda_n)^{n-2}, \\
&= (-1)^{n-1}(n-1)\,N(\mathfrak{a})^{-2}\,T^{n-2}
\end{aligned}
\]
und folglich nach (13) und (16)
\weq{17}{H = (-1)^{n-1}(n-1)\,T^{n-2}\Delta.}

\begin{center}
§ 88. \gesperrt{Idealklassen und Formenklassen.}
\end{center}

Wir haben in Bd. II, § 170 die \gesperrt{Äquivalenz} der Ideale und die \gesperrt{Idealklassen} erklärt. Danach waren zwei Ideale $\mathfrak{a}$, $\mathfrak{b}$ äquivalent. Wenn es eine ganze oder gebrochene \gesperrt{Zahl} des Körpers $\Omega$ gibt, die der Quotient von $\mathfrak{b}$ und $\mathfrak{a}$ ist, also:
\weq{1}{\eta\,\mathfrak{a} = \mathfrak{b},}
und es hat sich gezeigt, daß danach die Ideale in Klassen eingeteilt werden, und daß die \gesperrt{Zahl dieser Klassen endlich ist}. In manchen Fällen ist es zweckmäßig, den Äquivalenzbegriff etwas enger zu fassen und $\mathfrak{a}$ und $\mathfrak{b}$ nur dann äquivalent zu nennen, wenn es eine der Bedingung (1) genügende Zahl $\eta$ mit \gesperrt{positiver Norm} gibt. Dieser Unterschied kommt natürlich nicht in Betracht, wenn die Normen aller Zahlen positiv sind, wie z.\,B. im imaginären quadratischen Körper. Auch dann kommt er nicht in Betracht, wenn es \gesperrt{Einheiten mit der Norm} $-1$ gibt, weil, wenn $\varepsilon$ eine solche Einheit ist, entweder $N(\eta)$ oder $N(\varepsilon\eta)$ positiv ist und $\varepsilon$ und $\varepsilon\eta$ gleichzeitig der Bedingung (1) genügen.

Gibt es aber keine Einheiten, deren Norm $=-1$ ist, wohl aber andere Zahlen in $\Omega$ mit negativer Norm, so teilt sich bei der engeren Definition der Äquivalenz jede Idealklasse noch einmal in zwei Klassen. Wir wollen hier den engeren Äquivalenzbegriff festhalten, der z.\,B. bei manchen \gesperrt{reellen quadratischen Körpern} in Kraft tritt\footnote{Wenn nämlich die \gesperrt{Pell}sche Gleichung $t^2 - D\,u^2 = -4$ keine Lösung hat.}.

In (1) kann $\eta$ gebrochen sein. Ist aber $\alpha$ eine durch $\mathfrak{a}$ teilbare ganze Zahl, so ist $\eta\alpha$ eine durch $\mathfrak{b}$ \gesperrt{teilbare ganze Zahl}, und daraus ergibt sich: Wenn
\weq{2}{(\alpha_1,\ \alpha_2,\ \ldots,\ \alpha_n)}
eine Basis von $\mathfrak{a}$ ist, so ist
\[
(\beta_1,\ \beta_2,\ \ldots,\ \beta_n) = (\alpha_1\eta,\ \alpha_2\eta,\ \ldots,\ \alpha_n\eta)
\]
eine Basis von $\mathfrak{b}$, und wenn die erste Basis positiv ist, so ist es (bei dem engeren Äquivalenzbegriff) auch die zweite.
